\section{Arbitrary physical passage dimension when bypass is omitted}
\label{cp:section}

The global control theorem supplies a sharper obstruction to a power-only
channel proof. The following fields are responses of actual finite positive
dielectrics with physically normalized incoming waves. They satisfy positive
passage on arbitrarily large coherent subspaces. Their absolute bypass is
not uniformly controlled.

Here the experiment is deliberately narrower than the original channel
test: measure forward power on the square, but do not demand a bound on
the absolute flux elsewhere. The theorem is useful because it shows exactly
why the bypass condition cannot be dropped from a power-counting argument.
It does not establish the full counterexample on its own.

\begin{proposition}[Unbounded positive local-flux trace]
\label{cp:trace}
Fix $k>0$, $\varepsilon_+>1$, and the centered square $D$. In the
finite scalar dielectric class $1\leq\varepsilon\leq\varepsilon_+$,
the positive part of the section-flux operator has no uniform trace or
operator-norm bound in the physical incoming-power norm. This remains
true for radial dielectrics with a vacuum core containing $D$.
\end{proposition}
\begin{proof}
Let $R_D$ be the circumradius of $D$, choose odd $l$ with
$l(l+1)>k^2R_D^2$, and put $m=l+1$. Use real normalized axisymmetric
spherical harmonics and set
\[
 c_l=\frac{\partial_\theta Y_{l0}(\pi/2)}{\sqrt{l(l+1)}},
 \qquad d_m=Y_{m0}(\pi/2).
\]
These constants are nonzero. The radial positivity argument of
Theorem~\ref{cn:radial-bypass} gives $j_l(kr),j_m(kr)>0$ for
$0<r\leq R_D$. Hence
\begin{equation}\label{cp:positiveconstant}
 I_{lm}=\int_D\frac{j_l(kr)j_m(kr)}r\,dA>0,
 \qquad K_{lm}=4k\sqrt{m(m+1)}\,|c_ld_m|I_{lm}>0.
\end{equation}
The apparent singularity at the origin is integrable.

In the physical convention of Lemma~\ref{cg:vacuum-trace}, a free
unit-incoming-power TE wave, with its fixed phase absorbed into the
input basis, has electric field
$2k\sqrt{2Z}\,j_l(kr)X_{l0}$. Denote the actual core multipliers
by $g_l,g_m$ and the two physical incoming coefficients by $a,b$.
On the plane, the odd TE wave has azimuthal electric field and normal
magnetic field. The even TE wave has zero electric field and radial
tangential magnetic field there, with
\[
 H_r=i\frac{2k\sqrt{2Z}}Z\,b g_m\sqrt{m(m+1)}
             \frac{j_m(kr)}{kr}d_m.
\]
Since $\widehat\varphi\times\widehat r=-\widehat z$ on that plane,
the exact section form is
\begin{equation}\label{cp:two-mode-form}
 P_D(a,b)=-4k\sqrt{m(m+1)}c_ld_m I_{lm}
             \operatorname{Im}(a\overline b\,g_l\overline{g_m}).
\end{equation}
Its two eigenvalues are $\pm K_{lm}|g_lg_m|/2$.
Corollary~\ref{cc:core-gains} supplies actual finite radial devices
with $g_l,g_m$ approaching any common $G>0$, up to common phase.
The positive eigenvalue can therefore be arbitrarily large.
It bounds below the operator norm, the positive trace, and the
nuclear norm whenever those traces are defined; equivalently the
same divergence already holds on this two-dimensional compression.
\end{proof}

\begin{theorem}[Arbitrary coherent physical passage rank]
\label{cp:rank}
For every integer $N\geq1$ and $p_*>0$, there is a finite radial
dielectric with the same fixed $k$, $D$, and material bounds, and an
$N$-dimensional physical incoming space $W_N$, such that
\begin{equation}\label{cp:passage}
 \|Sa\|^2=\|a\|^2,\qquad P_D(a)\geq p_*\|a\|^2
       \quad(a\in W_N).
\end{equation}
Here $S$ is the complete scattering map. No common absolute-bypass
bound is asserted.
\end{theorem}
\begin{proof}
First construct local target fields. Choose distinct $q_j>k$ and
put $\alpha_j=\sqrt{q_j^2-k^2}$. In cylindrical coordinates define
\begin{equation}\label{cp:localfields}
 E_j=J_1(q_j\rho)
       \left[\cosh(\alpha_jz)+i\frac{k}{\alpha_j}
                          \sinh(\alpha_jz)\right]\widehat\varphi,
 \qquad H_j=\frac1{i\omega\mu_0}\operatorname{curl}E_j.
\end{equation}
The quotient $J_1(q_j\rho)/\rho$ is an entire function of $\rho^2$,
so the Cartesian vector field is smooth at the axis. Its divergence
vanishes, and the order-one Bessel equation gives
$(\Delta+k^2)E_j=0$. Both vacuum Maxwell curl equations follow.
At $z=0$, $\partial_zE_j=ikE_j$, whence every coherent combination
satisfies
\begin{equation}\label{cp:positivegram}
 S_z(a;\rho,0)=\frac1{2Z}
                 \left|\sum_{j=1}^Na_jJ_1(q_j\rho)\right|^2.
\end{equation}
Its section Gram matrix is positive definite. A combination
vanishing on $D$ vanishes on a radial interval; its Taylor series
then gives a Vandermonde system in the distinct $q_j^2$, after
extracting the nonzero factors $q_j$. All coefficients must vanish.

Next approximate these targets by a finite regular TE band. Each
axisymmetric azimuthal electric field in \eqref{cp:localfields} has
a regular TE spherical expansion with azimuthal order zero.
For the needed convergence, choose $R'>R_D$ avoiding zeros of the
finitely many possibly oscillatory $j_l(kR')$. All sufficiently
large orders are positive at $R'$. Expand the smooth tangential
trace on that sphere in $X_{l0}$. Angular projection of the vector
Helmholtz equation yields the regular radial factor $j_l(kr)$.
The smooth trace coefficients decay faster than any fixed power
of $l$; the large-order Bessel series gives a bound of the form
$C(r/R')^l$ on $j_l(kr)/j_l(kR')$, with polynomial factors for
first derivatives on any smaller ball. The expansion therefore
converges with its first derivatives on a ball containing $D$.
Projection of the radial equations and angular completeness
identify its sum with the original field.

Truncate all $N$ expansions at a common finite order $L$.
Uniform convergence of $E,H$ on $D$ preserves strict positivity
of their finite Hermitian Gram matrix once $L$ is sufficiently
large. In the physical incoming norm, the resulting coefficient
span $V_N$ consequently satisfies
\[
 \dim V_N=N,\qquad P_D^0(a)\geq p_N\|a\|^2,
 \qquad p_N>0.
\]
The constant $p_N$ may depend on $N$ and may be very small.
Apply Corollary~\ref{cc:core-gains} with a common real gain $G$
for every TE order through $L$. After common rephasing, the
actual finite-device section matrix converges to $G^2$ times
the free section matrix on $V_N$. Choose $G$ with
$G^2p_N>p_*$ and then a sufficiently accurate finite device.
This gives the second inequality of \eqref{cp:passage} on the
physical coefficient space $W_N=V_N$.
The device is finite, real, positive, and radial, so its complete
power-normalized scattering map is unitary. This gives the first
equality. The exponentially growing local targets themselves were
used only for compact approximation, never as physical inputs.
\end{proof}

\begin{remark}[Why this does not close the unrestricted channel target]
\label{cp:bypass}
For these finite angular-band devices, normal plane flux is absolutely
integrable. It is regular on bounded regions up to the finitely many
interfaces, and its exterior leading spherical Poynting vector is
radial and tangent to the plane, leaving an $O(r^{-3})$ normal tail.
Upper-half-space conservation with the physical radiation convention
gives
\[
 P_\Sigma=P_{\mathrm{out,upper}}-P_{\mathrm{in,upper}},
 \qquad |P_\Sigma(a)|\leq\|a\|^2.
\]
Consequently
\begin{equation}\label{cp:returncost}
 \int_{\Sigma\setminus D}|S_z(a)|
       \geq |P_D(a)-P_\Sigma(a)|
       \geq P_D(a)-\|a\|^2.
\end{equation}
The unbounded positive eigenvalue in Proposition~\ref{cp:trace}
therefore entails unbounded bypass. On a subspace satisfying a fixed
absolute-bypass bound, conservation only implies
$\|Q_D\|\leq1+\delta$. That operator-norm estimate alone cannot
bound the number of positive eigenvalues. The spatial information
in the absolute-bypass condition remains essential, and
Theorem~\ref{cb:finitecore} obstructs using the construction of
Theorem~\ref{cp:rank} at fixed bypass for arbitrarily large $N$.
\end{remark}
